\section*{Coda: the graph after gauge}
\addcontentsline{toc}{section}{Coda: the graph after gauge}

The object obtained in this chapter is an equivalence class made concrete by a declared reference,
then interpreted as a map between centered categorical fibers.  Raw coupling matrices and scalar
contact scores are respectively an input representation and a later compression of this object.
This is the first point at which edges from different constructions can enter a common coordinate
system while retaining their type.

Three distinctions should remain visible.  Gauge projection identifies the bivariate part of a
table; endpoint invariance is a separate response property.  Reference transport compares two
coordinate choices; population comparability depends on how the empirical references were
obtained.  Finally, an operator attached to an edge describes categorical action.  Static,
dynamic, probabilistic, computational, and physical readings require additional coordinates.

The physical reading is precise but conditional.  If $T_{ij}$ is a term in one compatible joint
Potts Hamiltonian, then adding $u(a)+v(b)+c$ redistributes energy between the pair table, the two
one-body fields, and the global zero without changing the Gibbs law.  The quotient therefore
isolates the part of an effective pair energy orthogonal to every reassignment into single-site terms.  It
generalizes the familiar freedom to choose an energy zero: the constant fixes the zero, while the
row and column terms fix how one-body and pair contributions are apportioned.

That statement concerns an effective Hamiltonian on sequence categories.  A microscopic-force
interpretation of $G_{ij}(a,b)$ additionally requires a declared
ensemble, a coarse-graining map from molecular configurations to residue identities, calibrated
energy units or a specified $k_{\mathrm B}T$, and evidence that the inferred sequence law is the
relevant equilibrium or evolutionary law.  In its native scope, $G_{ij}$ remains an exact
bivariate contrast operator: the part of the table invariant to one-body reparameterization, and
hence the appropriate object for comparison across models.

Those questions require a response experiment.  The next chapter therefore leaves the table as a
parameter and asks what happens when one residue is actually varied while another prediction is
observed.  In an additive Potts model the answer is fixed.  In a deep sequence model the answer can
move with the surrounding sequence, turning one edge operator into a field of operators over
backgrounds.
